\DocumentMetadata{
  pdfversion=2.0,
  pdfstandard=ua-2,
  lang=en-US,
  tagging=on,
  tagging-setup={math/setup=mathml-SE,tabsorder=structure}
}
\documentclass[11pt,letterpaper]{article}
\usepackage[margin=0.68in,includefoot]{geometry}
\usepackage{newcomputermodern}
\usepackage{amsmath,amscd,mathtools}
\usepackage{tikz,adjustbox}
\providecommand{\square}{\mdlgwhtsquare}
\usepackage[hidelinks]{hyperref}
\hypersetup{
  pdftitle={Probability PhD Exam, May 1998},
  pdfauthor={University of Florida Department of Mathematics},
  pdfsubject={Graduate mathematics examination},
  pdfdisplaydoctitle=true,
  bookmarksopen=true
}
\setcounter{secnumdepth}{0}
\setlength{\parindent}{0pt}
\setlength{\parskip}{0.28em}
\setlength{\emergencystretch}{3em}
\setlength{\leftmargini}{2.15em}
\setlength{\leftmarginii}{2.65em}
\setlength{\leftmarginiii}{3em}
\renewcommand{\labelenumi}{\textbf{\theenumi.}}
\pagestyle{empty}
\raggedbottom
\allowdisplaybreaks
\newenvironment{examproblems}[1]{%
  \begin{enumerate}%
  \setcounter{enumi}{#1}%
  \setlength{\topsep}{0.35em}%
  \setlength{\partopsep}{0pt}%
  \setlength{\itemsep}{0.48em}%
  \setlength{\parsep}{0.18em}%
}{\end{enumerate}}
\newenvironment{examsubparts}{%
  \begin{enumerate}%
  \setlength{\topsep}{0.18em}%
  \setlength{\partopsep}{0pt}%
  \setlength{\itemsep}{0.16em}%
  \setlength{\parsep}{0.1em}%
}{\end{enumerate}}
\newenvironment{examinstructions}{%
  \par\begingroup\itshape
}{%
  \par\endgroup\vspace{0.18em}
}
\makeatletter
\renewcommand\section{\@startsection{section}{1}{0pt}%
  {-1.25ex plus -.25ex minus -.1ex}%
  {0.55ex plus .1ex}%
  {\normalfont\large\bfseries}}
\renewcommand{\@maketitle}{%
  \newpage\null\vskip -1em%
  {\footnotesize\bfseries
    \href{https://www.ufl.edu/}{University of Florida}, \href{https://math.ufl.edu/}{Department of Mathematics}\par}%
  \vskip -0.5em%
  \tagstructbegin{tag=Title}%
    {\LARGE\bfseries\href{https://gma.math.ufl.edu/exams/probability-phd/}{Probability PhD Exam}, May 1998\par}%
  \tagstructend%
  \vskip 0.25em%
  \tagmcbegin{artifact}%
    \hrule height 1pt%
  \tagmcend%
  \vskip 1em%
}
\makeatother
\title{Probability PhD Exam, May 1998}
\author{}
\date{}
\begin{document}
\maketitle
\thispagestyle{empty}
\begin{examproblems}{0}
\item
Let~\(X\) and~\(Y\) be independent and uniform on \([0,1]\). Find:
\begin{examsubparts}
\item[\textnormal{a)}]
the distribution of \(X+Y\).
\item[\textnormal{b)}]
the conditional density of~\(X\) given \(X+Y=z\).
\end{examsubparts}
\item
Let \(X_0,X_1,\ldots,X_n,\ldots\) be i.i.d. with mean~\(m\). Let~\(N\) be Poisson with mean~\(\lambda\), independent of the~\(X\)’s. Define 
\[
Y=X_0+\cdots+X_N.
\]
 Prove that~\(Y\) is integrable. Find \(E(Y)\).
\item
State precisely (all for i.i.d.):
\begin{examsubparts}
\item[\textnormal{a)}]
the strong law of large numbers and prove it.
\item[\textnormal{b)}]
the central limit theorem.
\item[\textnormal{c)}]
the law of the iterated logarithm.
\end{examsubparts}
\item
Define a martingale relative to a filtration \(\{\mathcal F_n\}\), \(n=0,1,2,\ldots\). Define a stopping time relative to this filtration. State the optional sampling theorem. State and prove the martingale convergence theorem.
\item
What is an infinitely divisible distribution on \(\mathbb R^1\). Prove that a non-trivial infinitely divisible distribution cannot have compact support.
\item
Define a standard Brownian Motion~\(W\) starting at 0. Define new processes \(W_1,W_2,W_3\) by: 
\[
W_1(t)=cW\!\left(\frac{t}{c^2}\right),\qquad t\geq0,\quad c\text{ real},\ c\neq0,
\]
 
\[
W_2(t)=tW\!\left(\frac1t\right),\qquad t>0,\quad W_2(0)=0,
\]
 
\[
W_3(t)=
\begin{cases}
W(1)-W(1-t),&0\leq t\leq1,\\
W(t),&\text{otherwise.}
\end{cases}
\]
 Show that \(W_1,W_2,W_3\) are the standard Brownian motions.
\item
Let \(\varphi(t)=\int e^{itx}\,\mu(dx)\) be a characteristic function where~\(\mu\) is a probability measure on \(\mathbb R^1\).
\begin{examsubparts}
\item[\textnormal{a)}]
Suppose \(|\varphi(t)|=1\) for some \(t\neq0\). Then prove that unless \(|\varphi(t)|\equiv1\), there is a smallest \(t_0>0\) and a~\(d\) such that~\(\mu\) is concentrated on the set \(\left\{d+\frac{2\pi j}{t_0}\right\}\), \(j=0,\pm1,\pm2,\ldots\).
\item[\textnormal{b)}]
Using a) prove that \(|\cos t|\) is not a characteristic function. Hint: \(\cos t\) and \(\cos^2t\) are characteristic functions.
\end{examsubparts}
\end{examproblems}
\end{document}
